\section{Why equation induction cannot replace numerical cases}
\label{sec:product-obstruction}

A tempting simplification changes the case introduction from a proof of
$B(t,Sz)$ under $B(t,0)$ to a proof of $B(t,z)$ under $B(t,0)$, with $z$
fresh.  The modified rule is sound for the ordinary conditional-closure
reading without any numerical coverage axiom.  It is nevertheless strictly
weaker when combined with the remaining native rules and literal-equation
induction.  This section gives both a semantic obstruction and a closed
native formula separating the two calculi.  These results do not assert
that every possible alternative interpretation of the weaker theory is
impossible.

\subsection{A product model of all the equational arithmetic}

Let $\mathcal A$ have numerical domain $\mathbb N\times\mathbb N$.  Interpret
zero, successor, and iteration coordinatewise:
\[
 0^{\mathcal A}=(0,0),\qquad
 S^{\mathcal A}(a,b)=(a+1,b+1).
\]
For each finite iterator body, the first output coordinate is the usual
natural-number iteration of its first-coordinate interpretation, and the
second coordinate is defined in the same way.  Equality is actual equality
of ordered pairs.  Use ordinary two-valued first-order semantics; a
classical structure is also a sound test of intuitionistic derivability.

More explicitly, if $\rho$ assigns pairs to the free variables, write
$\rho_i$ for its $i$th-coordinate assignment.  Structural induction on a
scalar term $t$ gives
\begin{equation}
 \llbracket t\rrbracket_{\mathcal A,\rho}
 =\bigl(\llbracket t\rrbracket_{\mathbb N,\rho_1},
         \llbracket t\rrbracket_{\mathbb N,\rho_2}\bigr).
 \label{eq:product-term-factorization}
\end{equation}
The iterator case follows by the defining coordinatewise interpretation,
including the coordinatewise assignment to its bound state variable.
Consequently, for every literal equation $E(n,\vec p)$ and every fixed
parameter assignment, there are unary predicates $E_1,E_2$ on the standard
natural numbers such that
\[
 \mathcal A\models E((a,b),\vec p)
 \quad\Longleftrightarrow\quad E_1(a)\ \text{and}\ E_2(b).
\]
This assertion concerns equations.  Arbitrary formulas containing
implication or quantifiers need not factor this way.

Both iterator computation equations hold in $\mathcal A$, successor is
injective, and $0\ne S0$.  Moreover, every literal-equation induction
instance is valid.  To prove the latter claim, suppose its base and step
hold.  The base supplies $E_1(0)$ and $E_2(0)$.  Given $E_1(a)$, evaluate
the product step at $(a,0)$, using $E_2(0)$, to obtain $E_1(a+1)$.
Similarly, evaluating the step at $(0,b)$ gives
$E_2(b)\Rightarrow E_2(b+1)$.  Ordinary induction in each standard
coordinate yields $E_1(a)$ and $E_2(b)$ for all $a,b$.  This proves the
product induction conclusion.

The argument also proves soundness of contextual equation induction.
Fix an assignment making every retained assumption in $\Gamma$ true.
Freshness of the induction variable means that changing its pair value
does not change the truth of $\Gamma$.  The preceding argument therefore
applies to the base and step under that same assignment.  No restriction to
equational side assumptions is required.

Numerical coverage fails.  The element $(0,1)$ is neither $(0,0)$ nor a
successor: every successor has two positive coordinates.  Thus
\begin{equation}
 \mathcal A\not\models
 \forall x\bigl(x=0\lor\exists y\,x=Sy\bigr).
 \label{eq:product-no-coverage}
\end{equation}
Adding all valid predecessor computation equations does not remove this
obstruction: coordinatewise predecessor already satisfies
$P(0)=0$ and $P(Sx)=x$ in this model.

\subsection{The failure is visible inside the native language}

Interpret the native binders by their ordinary readings
\[
 Q_yB:=B(0)\to\forall y\,B(y),\qquad
 D_{x,y}B:=\exists x\bigl(B(x,0)\to\forall y\,B(x,y)\bigr).
\]
Hypothesis, equality transport, QE, and DE are sound in $\mathcal A$.
The two computations and equation progress are sound by the preceding
argument.  The proposed direct-target DI rule is sound as well: its fresh
variable already ranges over the whole domain.  Hence every theorem of
that modified calculus is valid in $\mathcal A$.

Write $1=S0$ and $C(v,a,b)=I(v,a;u.b)$, with $u$ absent from $b$.
Define the following native formulas, where the displayed $E$ is a single
literal equation:
\begin{align*}
 E(n,y,p)&:\quad
 C\bigl(p,C(n,y,1),C(n,1,0)\bigr)=0,\\
 P(n,p)&:=Q_yE(n,y,p),\\
 G(p)&:=Q_nP(n,p),\\
 F&:=Q_pG(p).
\end{align*}
The last formula is closed.

The original calculus proves $F$ without equation induction.  At a
successor $n=Sz$, computation makes $E(Sz,y,p)$ equivalent to
$C(p,1,0)=0$, independently of $y$.  Thus $P(Sz,p)$ follows by
conditional-closure introduction from that same equation as a retained
assumption.  The original case introduction, with an unused witness,
therefore proves $Q_nP(n,p)$: its successor premise is already proved and
its retained zero premise can be ignored.  A further such introduction,
now in the parameter $p$, proves $F$.  Only case DI, DE, QE, computations,
and equality reasoning are needed for these derived closure introductions.

In $\mathcal A$, take $p_*=(0,1)$.  Direct coordinatewise computation gives
\[
 \mathcal A\models E(n,y,p_*)
 \quad\Longleftrightarrow\quad
 n_1=0,\quad n_2>0,\quad y_1=0.
\]
Consequently $P(0,p_*)$ is true, because its antecedent is false.
Likewise $P(Sz,p_*)$ is true for every $z$.  But $P((0,1),p_*)$ is
false: its antecedent holds at $y=0$, whereas its universally quantified
consequent fails at $y=(1,0)$.  Thus $G(p_*)$ is false.

Finally, $G(0)$ is true.  At $p=0$ and $n=0$, the equation $E$ becomes
$y=0$, so $P(0,0)=Q_y(y=0)$ is false in this nontrivial structure.  That
false antecedent makes $G(0)$ true.  Since $G(p_*)$ is false, the closed
formula $F=Q_pG$ is false.  Soundness of the modified calculus in
$\mathcal A$ therefore proves that it cannot derive $F$, while the
original calculus does derive $F$.

\subsection{What remains possible}

The obstruction distinguishes algebraic Boolean identities from global
logical alternatives.  For example, equation induction proves
$t(C(n,a,b))=C(n,t(a),t(b))$ for each fixed term context $t$.
It can also prove that the Boolean selector value $b=C(n,0,1)$ satisfies
$C(b,0,1)=b$.  Such identities remain valid for the mixed Boolean pairs
$(0,1)$ and $(1,0)$.  They do not entail $b=0\lor b=1$.
Accordingly, improving the equational predecessor construction cannot by
itself recover the missing numerical case rule.

There is a separate, unfinished route through recursion uniqueness.
Much of the scalar predecessor bootstrap can be recast as uniqueness of
an explicitly displayed recurrence. Write
$Z(u)=C(u,1,0)$, $N(n)=I(n,0;u.Z(u))$, and
$a+b=I(b,a;u.Su)$, with $2a=a+a$. For example, if
$j(u)=Z(u)+2N(u)$, $M(n)=I(n,0;j)$, and
$G(n)=I(n,0;u.u+M(u)+1)$, the difficult comparison
$M(G(n))=N(n)$ can compare two recurrences with the common update
$k(v)=I(Sv,v;j)$.  Finite Boolean computations give
$k(N(n))=Z(N(n))$.  This is useful progress on that route, but it does not
complete equality reflection or establish equivalence of a new
uniqueness-only theory with iPRA; no such conclusion is claimed here.
